% layout-demo.tex
% Compile with: pdflatex layout-demo.tex
% Purpose: compare 3 formula+code layout options before writing the full derivation doc.
% All examples use the same real formula (physical state transition) and real code lines.

\documentclass[11pt,a4paper]{article}

\usepackage{amsmath, amssymb}
\usepackage{xcolor}
\usepackage{listings}
\usepackage{tcolorbox}
\tcbuselibrary{listings, skins, breakable}
\usepackage{geometry}
\geometry{margin=2.5cm}
\usepackage{parskip}
\usepackage{booktabs}
\usepackage{caption}

% ── Python listing style ──────────────────────────────────────────────────────
\definecolor{codebg}{RGB}{245,245,245}
\definecolor{codeframe}{RGB}{200,200,200}
\definecolor{keywordcolor}{RGB}{0,0,180}
\definecolor{commentcolor}{RGB}{80,120,80}
\definecolor{stringcolor}{RGB}{160,0,0}

\lstdefinestyle{python}{
  language=Python,
  basicstyle=\ttfamily\small,
  keywordstyle=\color{keywordcolor}\bfseries,
  commentstyle=\color{commentcolor}\itshape,
  stringstyle=\color{stringcolor},
  showstringspaces=false,
  breaklines=true,
  frame=none,
  backgroundcolor=\color{codebg},
  numbers=left,
  numberstyle=\tiny\color{gray},
  numbersep=6pt,
  xleftmargin=12pt,
}

% ── tcolorbox styles for Option B ─────────────────────────────────────────────
\tcbset{
  formulabox/.style={
    enhanced,
    colback=white,
    colframe=black!50,
    arc=2pt,
    boxrule=0.5pt,
    left=6pt, right=6pt, top=4pt, bottom=4pt,
  },
  codebox/.style={
    enhanced,
    colback=codebg,
    colframe=black!50,
    arc=2pt,
    boxrule=0.5pt,
    top=0pt, bottom=0pt, left=0pt, right=0pt,
    listing only,
    listing options={style=python, numbers=left, xleftmargin=12pt},
  },
  pairedbox/.style={
    enhanced,
    breakable,
    colback=white,
    colframe=black!60,
    arc=3pt,
    boxrule=0.6pt,
    fonttitle=\small\ttfamily,
    attach boxed title to top left={yshift=-2mm, xshift=4mm},
    boxed title style={colback=black!70, colframe=black!70, arc=2pt, boxrule=0pt,
                        left=4pt, right=4pt},
    coltitle=white,
  },
}

% ── Convenience command: inline filename ──────────────────────────────────────
\newcommand{\file}[1]{\texttt{\small #1}}

% =============================================================================
\begin{document}

\title{\textbf{Layout Demo} \\[4pt]
       \large Formula + Code Pairing: Three Options}
\author{}
\date{}
\maketitle

\noindent
Each option shows the \textit{same} formula and \textit{same} code lines so the visual
difference is the only variable. The formula is the physical state transition in the
augmented model; the code is from \file{gantry\_interconnect\_dynamic.py}.

\bigskip
\hrule
\bigskip

% =============================================================================
\section*{Option A \quad Equation + shaded code block, stacked}

Plain equation environment, then an \texttt{lstlisting} block with a shaded background
and a title line showing filename and line numbers. Minimal packages (only
\texttt{listings} + \texttt{xcolor}).

% ── Formula ──
\begin{equation}
  \hat{x}^{\mathrm{phys}}_{k+1}
  = f_{\mathrm{LPV}}\!\bigl(
      S_{\mathrm{phys}}\,\hat{x}_k,\;
      \hat{u}_k
    \bigr)
  \label{eq:phys_transition}
\end{equation}

where $S_{\mathrm{phys}} = \texttt{selection\_matrix(PHY\_IX, nxd)}$ selects the six
physical-state rows of $\hat{x}_k$, and the output $\hat{x}^{\mathrm{phys}}_{k+1}$ is
written back into those rows via $E_{\mathrm{phys}} = \texttt{expansion\_matrix(PHY\_IX, nxd)}$.

% ── Code ──
\begin{lstlisting}[
  style=python,
  backgroundcolor=\color{codebg},
  frame=single,
  framerule=0.4pt,
  rulecolor=\color{codeframe},
  title={\small\ttfamily gantry\_interconnect\_dynamic.py \quad lines 267--269},
  captionpos=t,
  firstnumber=267,
]
ic.connect_signals("x", phy_block, "concat", selection_matrix(PHY_IX, nxd))
ic.connect_block_signals(phy_block, ["u"], [])
ic.connect_signals(phy_block, "xp", "additive", expansion_matrix(PHY_IX, nxd))
\end{lstlisting}

\bigskip
\hrule
\bigskip

% =============================================================================
\section*{Option B \quad \texttt{tcolorbox} split panel}

A two-part colored box: white top panel for the formula, gray bottom panel for the code.
The title on the right identifies the source file. Requires \texttt{tcolorbox}.

% ── Paired box ──
\begin{tcolorbox}[
  pairedbox,
  title={gantry\_interconnect\_dynamic.py \quad lines 267--269},
]
% top: formula
\vspace{2pt}
\begin{equation*}
  \hat{x}^{\mathrm{phys}}_{k+1}
  = f_{\mathrm{LPV}}\!\bigl(
      S_{\mathrm{phys}}\,\hat{x}_k,\;
      \hat{u}_k
    \bigr)
\end{equation*}
\vspace{-4pt}

\noindent\small
$S_{\mathrm{phys}}$: \texttt{selection\_matrix(PHY\_IX, nxd)}
\hfill
$E_{\mathrm{phys}}$: \texttt{expansion\_matrix(PHY\_IX, nxd)}

\tcblower % ── bottom: code ──
\begin{lstlisting}[style=python, firstnumber=267, numbers=left, xleftmargin=12pt]
ic.connect_signals("x", phy_block, "concat", selection_matrix(PHY_IX, nxd))
ic.connect_block_signals(phy_block, ["u"], [])
ic.connect_signals(phy_block, "xp", "additive", expansion_matrix(PHY_IX, nxd))
\end{lstlisting}
\end{tcolorbox}

\bigskip
\hrule
\bigskip

% =============================================================================
\section*{Option C \quad Numbered equation with matching implementation label}

Standard numbered equation, then a titled code block whose number echoes the equation.
Works well with cross-referencing (\texttt{impl.~\ref{eq:phys_transition_c}}).

% ── Formula ──
\begin{equation}
  \hat{x}^{\mathrm{phys}}_{k+1}
  = f_{\mathrm{LPV}}\!\bigl(
      S_{\mathrm{phys}}\,\hat{x}_k,\;
      \hat{u}_k
    \bigr)
  \label{eq:phys_transition_c}
\end{equation}

% ── Code ──
\begin{lstlisting}[
  style=python,
  backgroundcolor=\color{codebg},
  frame=single,
  framerule=0.4pt,
  rulecolor=\color{codeframe},
  title={%
    \small\textbf{impl.\,\ref*{eq:phys_transition_c}}\quad
    \ttfamily gantry\_interconnect\_dynamic.py \quad lines 267--269},
  captionpos=t,
  firstnumber=267,
]
ic.connect_signals("x", phy_block, "concat", selection_matrix(PHY_IX, nxd))
ic.connect_block_signals(phy_block, ["u"], [])
ic.connect_signals(phy_block, "xp", "additive", expansion_matrix(PHY_IX, nxd))
\end{lstlisting}

\bigskip
\hrule
\bigskip

% =============================================================================
\section*{Candidate sections for the full document}

Once you pick a layout style, the full derivation doc will contain these sections
in this order:

\begin{enumerate}
  \item Augmented state partitioning $\hat{x}_k = [x^\mathrm{phys}_k;\, x^\mathrm{aug}_k]$
  \item State transition -- physical block (LPV gantry)
  \item State transition -- augmented block (ANN)
  \item Output equation (two-block additive)
  \item Normalisation ($C_{d,\mathrm{norm}}$, $\bar{u}$, $\bar{y}$)
  \item Encoder initialisation (reconstructability map, Eqs.\ 16--17)
  \item Training objective (truncated multi-step loss)
\end{enumerate}

\end{document}
